\chapter{Editing a schedule by hand}\label{ch:editing}

Planners know things the solver does not: a machine that is about to be serviced, a customer who called, a
team that would rather work one long run than two short ones. PlanWise lets you change any option directly.
\textbf{You decide the machine and the order of runs; the computer works out the times}, re-checks the whole
schedule and scores it, after every single change. The original option is never modified: your edit is saved
as a new option.

\section{Starting the editor}

On \ui{Schedule \& compare}, pick the option to start from and click \ui{Edit this schedule by hand} (or
\ui{Edit by hand} in an order's explanation, or \ui{Edit} in a quick look). Editing is available for
five-stage schedules of open periods.

\shot{editor_start}{The editor before any change: the toolbar, the figures ``as original'' and the Gantt chart.}{editor_start}

The editor shows a sticky toolbar, the five headline figures with their difference from the original, the
Gantt chart and, on the right, how to edit (until you select a run).

\section{Changing runs}\label{sec:edit-runs}

There are two ways to change a run:

\begin{itemize}
  \item \textbf{Drag} a bar to another machine of the same stage or to another time (in 15-minute steps). If
        the target machine cannot make that item, the bar turns red and cannot be dropped there.
  \item \textbf{Select} a bar (click it, or move to it with the arrow keys and press \kbd{Enter}) and use the
        run panel:
\end{itemize}

\shotf[0.72\linewidth]{editor_runpanel}{The run panel of a machining run moved from CNC\_1 to CNC\_2.}{editor_runpanel}

\begin{description}[style=nextline, leftmargin=1.2em, font=\bfseries\color{brand}]
  \item[Lock this run (L)] Keeps the run exactly as it is when you later use \ui{Optimise the rest}.
  \item[Original] Where the run was in the original option.
  \item[Machine] Move the run to another eligible machine. If only one machine can make it, the field says so.
  \item[Order on the machine] \ui{Earlier} / \ui{Later} swap it with the neighbouring run on its machine
        (``Run 5 of 29 on CNC\_2'').
  \item[Not before] An earliest start: the run starts at the earliest valid slot from that time. \ui{Not
        earlier than now} pins it to its current start; \ui{Clear} removes the constraint.
\end{description}

\section{Seeing the effect immediately}

\shot{editor}{After moving one run: the edit is valid, 60 cost units cheaper but 2.9 hours more late in total.}{editor}

\begin{callouts}
  \item[1] \textbf{Undo, redo, reset} (\kbd{Ctrl}\,\kbd{Z}, \kbd{Ctrl}\,\kbd{Shift}\,\kbd{Z}). Every change can be
           undone.
  \item[2] \textbf{How times are rebuilt}: \ui{Keep other times} (untouched runs keep their time; only runs that
           get pushed move) or \ui{Pack early} (every run starts as early as the current order allows).
  \item[3] \textbf{Exit}: leave the editor; unsaved changes are discarded after a confirmation.
  \item[4] \textbf{Optimise the rest\dots}: let CP-SAT re-sequence everything you did not lock.
  \item[5] \textbf{Save as new option}.
  \item[6] \textbf{Impact}: each figure with its change versus the original (green better, amber worse). The
           toolbar repeats it in one line, for example ``Valid · Late +2.9 h · Cost −60''.
  \item[7] \textbf{Run panel} of the selected run.
\end{callouts}

After every change the backend rebuilds the times (respecting shifts, setups, mold cycles and maintenance,
semi-finished stock and warehouse caps), runs the independent validator and scores the result. Changed runs
get a dashed outline, and a \ui{Changes} panel lists your edits and every run that moved as a consequence. A
banner explains shift effects (``The edit closes 1 shift(s): CNC\_1 16:00--00:00 Mon, 11/23. Operating cost
down 1,446.'').

If the edit breaks a rule, it cannot be saved and the problem is explained:

\begin{center}
\begin{tabularx}{\linewidth}{>{\bfseries}p{3.4cm}X}
\toprule
\rowcolor{white}Problem & Example message \\
\midrule
Warehouse full & Finished stock of F silver exceeds the warehouse cap at \dots \\
Material missing & Semi-finished F\_HEAT runs out at \dots: a run uses it before it is made. \\
Buffer full & Semi-finished buffer F\_CAST exceeds its capacity at \dots \\
Cannot ship & Not enough R black to ship at \dots \\
Beyond the period & Run \dots\ is pushed past the end of the period. \\
No slot & Cannot place \dots: they wait for semi-finished stock not made before them, or the period is full. Try moving the upstream run earlier. \\
\bottomrule
\end{tabularx}
\end{center}

Click \ui{Show} next to a problem to select the run concerned, then undo or keep editing until the schedule is
valid again.

\section{Optimise the rest}\label{sec:optimise-rest}

Often you only care about a few runs. Lock them (or simply edit them) and let CP-SAT re-sequence everything else:

\shotf[0.55\linewidth]{optimize_dialog}{Optimise the rest: 1 run kept, 296 runs CP-SAT may change.}{optimize_dialog}

\begin{description}[style=nextline, leftmargin=1.2em, font=\bfseries\color{brand}]
  \item[Kept / CP-SAT may change] How many runs stay fixed and how many may move.
  \item[Keep the runs I edited] On by default; locked runs are always kept.
  \item[Frozen zone] Runs starting before this point stay: \ui{None}, \ui{Until now} (while the period is
        running), \ui{First 24 h}, \ui{First 48 h} or \ui{Until\dots} a chosen time. Use it to leave the shifts
        already handed to the floor untouched.
  \item[Time allowed] 30~s, 2~min or 5~min. Longer usually finds more; you can keep working meanwhile.
  \item[Short name] Optional name of the resulting option.
  \item[Also save the current edit as its own option] To compare your hand edit with the optimised result.
\end{description}

The optimisation runs in the background like any solver run. Its result is a new option that is \emph{never
worse} than your edit; when it finishes a message says whether something better was found. With spare
capacity, CP-SAT may add build-ahead runs for long-dated orders to reduce safety-stock shortfall.

\section{Saving the edit}

\shotf[0.5\linewidth]{save_option_dialog}{Save as a new option.}{save_option_dialog}

Give the edit a short name (``Machining run moved to CNC lathe 2'') and, optionally, why you made it. The
original option is kept. The edit is ranked with the other options under the same policy and can be compared
and committed as usual; its card says \ui{Manual edit} and from which option it was made. A hand edit that
turns out to be the best schedule is recommended like any other --- in the example the edit became the
``Lowest cost'' option and was committed.
